% ====================================================================
% Khung de sinh tu bang dac ta: 3_Bang_dac_ta.docx
% Moi cau da co san Chu de / Don vi kien thuc / Yeu cau can dat / Muc do.
% Viec can lam: (1) dien noi dung cau hoi va loi giai vao cac cho TODO,
% (2) tra MAPID de thay dau '?' trong ID6.
% ====================================================================
\documentclass[FileMain.tex]{subfiles}
\gdef\sophong{Sở GD \& ĐT Gia Lai}
\gdef\truong{Trường THPT Chi Lăng}
\gdef\truongh{Trường Mầm non, THCS, THPT Sao Việt}
\gdef\monhoc{Khoa học tự nhiên 6} % ví dụ: Khoa học tự nhiên 6
\gdef\ngaykt{20/10/2026} % ví dụ: 04/02/2026
\gdef\nh{2025 - 2026} % ví dụ: 2025 - 2026
\gdef\thoigian{45}% ví dụ: 45
\gdef\made{101} % ví dụ: 758
\setcounter{section}{0}% ví dụ đề 1 là 0, đề 2 là 1.
%\tatloigiai
%\hienthiloigiai
%\dongkeloigiai
\begin{document}
\section[Kiểm tra định kì giữa học kì I Khoa học tự nhiên 6 - Mã đề \made]{Kiểm tra định kì giữa học kì I} % ví dụ Kiểm tra giữa kỳ 1

%%%==============Phần trắc nghiệm nhiều lựa chọn==============%%%
\subsection{Bài tập trắc nghiệm nhiều lựa chọn}\textit{\large Thí sinh trả lời từ câu 1 đến 12. Mỗi câu thí sinh chỉ chọn một phương án}
\Opensolutionfile{ansex}[Ans/LGEX-KTDK_KHTN6_MoDau_DaDangChat_MADE101]
\Opensolutionfile{ans}[Ans/Ans-KTDK_KHTN6_MoDau_DaDangChat_MADE101]
%%%=========EX_1================%%%
% Chu de : Mở đầu về Khoa học tự nhiên
% Don vi : 1.1. Giới thiệu về Khoa học tự nhiên. Các lĩnh vực chủ yếu của KHTN. Vật sống – vật không sống
% Muc do : Nhận biết
% YCCD : – Nêu được khái niệm khoa học tự nhiên; nhận ra được đối tượng nghiên cứu của khoa học tự nhiên.
\begin{ex}%[6KCN1-1]
Khoa học tự nhiên nghiên cứu về đối tượng nào sau đây?
\choice
{\True Các sự vật, hiện tượng và quy luật của thế giới tự nhiên}
{Sự hình thành và phát triển của xã hội loài người}
{Các quy tắc ứng xử giữa người với người}
{Các tác phẩm hội hoạ và âm nhạc}
\loigiai{
Khoa học tự nhiên là ngành khoa học nghiên cứu các sự vật, hiện tượng, quy luật của tự nhiên cùng ảnh hưởng của chúng đến cuộc sống con người và môi trường.
}
\end{ex}

%%%=========EX_2================%%%
% Chu de : Mở đầu về Khoa học tự nhiên
% Don vi : 1.1. Giới thiệu về Khoa học tự nhiên. Các lĩnh vực chủ yếu của KHTN. Vật sống – vật không sống
% Muc do : Nhận biết
% YCCD : – Trình bày được các lĩnh vực chủ yếu của khoa học tự nhiên: Vật lí, Hoá học, Sinh học, Thiên văn học, Khoa học Trái Đất.
\begin{ex}%[6KCN1-2]
Lĩnh vực nào sau đây KHÔNG thuộc khoa học tự nhiên?
\choice
{Sinh học}
{Thiên văn học}
{\True Lịch sử}
{Hoá học}
\loigiai{
Năm lĩnh vực chủ yếu của khoa học tự nhiên là Vật lí, Hoá học, Sinh học, Thiên văn học và Khoa học Trái Đất. Lịch sử thuộc nhóm khoa học xã hội.
}
\end{ex}

%%%=========EX_3================%%%
% Chu de : Mở đầu về Khoa học tự nhiên
% Don vi : 1.1. Giới thiệu về Khoa học tự nhiên. Các lĩnh vực chủ yếu của KHTN. Vật sống – vật không sống
% Muc do : Thông hiểu
% YCCD : – Phân biệt được vật sống và vật không sống dựa vào các đặc điểm đặc trưng (trao đổi chất, lớn lên, sinh sản, cảm ứng); phân biệt được đối tượng của KHTN với các ngành khoa học khác.
\begin{ex}%[6KCH1-3]
Vật nào sau đây là vật sống?
\choice
{Robot lau nhà}
{\True Cây cà chua}
{Hòn đá cuội}
{Chiếc xe đạp}
\loigiai{
Cây cà chua có đầy đủ các đặc điểm của vật sống: trao đổi chất, lớn lên, sinh sản và cảm ứng. Robot chuyển động được nhưng không trao đổi chất và không sinh sản.
}
\end{ex}

%%%=========EX_4================%%%
% Chu de : Mở đầu về Khoa học tự nhiên
% Don vi : 1.2. An toàn trong phòng thực hành. Sử dụng dụng cụ đo, kính lúp và kính hiển vi quang học
% Muc do : Nhận biết
% YCCD : – Trình bày được công dụng và cách sử dụng một số dụng cụ đo thông dụng (thước, bình chia độ, cân, nhiệt kế), kính lúp và kính hiển vi quang học.
\begin{ex}%[6KCN2-2]
Dụng cụ nào sau đây được dùng để đo thể tích chất lỏng trong phòng thực hành?
\choice
{Cân điện tử}
{\True Bình chia độ}
{Nhiệt kế}
{Lực kế}
\loigiai{
Bình chia độ (ống đong) dùng để đo thể tích chất lỏng; cân đo khối lượng, nhiệt kế đo nhiệt độ, lực kế đo lực.
}
\end{ex}

%%%=========EX_5================%%%
% Chu de : Mở đầu về Khoa học tự nhiên
% Don vi : 1.2. An toàn trong phòng thực hành. Sử dụng dụng cụ đo, kính lúp và kính hiển vi quang học
% Muc do : Nhận biết
% YCCD : – Trình bày được công dụng và cách sử dụng một số dụng cụ đo thông dụng (thước, bình chia độ, cân, nhiệt kế), kính lúp và kính hiển vi quang học.
\begin{ex}%[6KCN2-2]
Kính hiển vi quang học được dùng để quan sát
\choice
{các thiên thể ở rất xa trên bầu trời}
{các chi tiết nhỏ của con tem, gân lá bằng cách phóng to vài chục lần}
{\True các vật rất nhỏ mà mắt thường không nhìn thấy được như tế bào, vi khuẩn}
{kích thước thật của các vật trong đời sống hằng ngày}
\loigiai{
Kính hiển vi quang học phóng đại hàng trăm đến hàng nghìn lần, dùng quan sát tế bào, vi khuẩn. Quan sát gân lá, con tem là công dụng của kính lúp; quan sát thiên thể dùng kính thiên văn.
}
\end{ex}

%%%=========EX_6================%%%
% Chu de : Mở đầu về Khoa học tự nhiên
% Don vi : 1.2. An toàn trong phòng thực hành. Sử dụng dụng cụ đo, kính lúp và kính hiển vi quang học
% Muc do : Thông hiểu
% YCCD : – Xác định được giới hạn đo (GHĐ) và độ chia nhỏ nhất (ĐCNN) của dụng cụ đo; lựa chọn được dụng cụ đo phù hợp với phép đo.
\begin{ex}%[6KCH2-3]
Trên một chiếc cân đồng hồ có ghi 5 kg -- 20 g. Số 20 g cho biết điều gì?
\choice
{Giới hạn đo của cân}
{Khối lượng của đĩa cân}
{\True Độ chia nhỏ nhất của cân}
{Khối lượng lớn nhất mỗi lần cân thêm được}
\loigiai{
Số lớn (5 kg) là giới hạn đo -- khối lượng lớn nhất cân đo được; số nhỏ (20 g) là độ chia nhỏ nhất -- giá trị giữa hai vạch chia liên tiếp.
}
\end{ex}

%%%=========EX_7================%%%
% Chu de : Đa dạng chất. Các thể của chất
% Don vi : 2.1. Sự đa dạng của chất. Ba thể cơ bản của chất và đặc điểm
% Muc do : Nhận biết
% YCCD : – Nêu được sự đa dạng của chất; chất có ở xung quanh ta, trong mọi vật thể tự nhiên và vật thể nhân tạo.
\begin{ex}%[6K1N1-1]
Vật thể nào sau đây là vật thể tự nhiên?
\choice
{\True Cây mía}
{Chiếc bàn gỗ}
{Cốc thuỷ tinh}
{Chiếc áo len}
\loigiai{
Cây mía tồn tại sẵn trong tự nhiên. Bàn gỗ, cốc thuỷ tinh, áo len đều do con người chế tạo nên là vật thể nhân tạo.
}
\end{ex}

%%%=========EX_8================%%%
% Chu de : Đa dạng chất. Các thể của chất
% Don vi : 2.1. Sự đa dạng của chất. Ba thể cơ bản của chất và đặc điểm
% Muc do : Nhận biết
% YCCD : – Nêu được ba thể (trạng thái) cơ bản của chất: rắn, lỏng, khí kèm ví dụ minh hoạ ở điều kiện thường.
\begin{ex}%[6K1N1-2]
Ở điều kiện thường, chất nào sau đây tồn tại ở thể khí?
\choice
{Đường ăn}
{\True Oxygen}
{Nước}
{Sắt}
\loigiai{
Ở điều kiện thường oxygen là chất khí; đường ăn và sắt ở thể rắn, nước ở thể lỏng.
}
\end{ex}

%%%=========EX_9================%%%
% Chu de : Đa dạng chất. Các thể của chất
% Don vi : 2.1. Sự đa dạng của chất. Ba thể cơ bản của chất và đặc điểm
% Muc do : Thông hiểu
% YCCD : – Trình bày và so sánh được đặc điểm về hình dạng, thể tích, khả năng bị nén và khả năng lan truyền của chất ở thể rắn, lỏng, khí.
\begin{ex}%[6K1H1-3]
Đặc điểm nào sau đây là của chất ở thể lỏng?
\choice
{Có hình dạng và thể tích xác định}
{\True Có thể tích xác định, hình dạng thay đổi theo vật chứa}
{Không có hình dạng và thể tích xác định}
{Chiếm toàn bộ thể tích của vật chứa và rất dễ bị nén}
\loigiai{
Chất lỏng có thể tích xác định nhưng hình dạng phụ thuộc vật chứa. Phương án A là thể rắn; C và D là thể khí.
}
\end{ex}

%%%=========EX_10================%%%
% Chu de : Đa dạng chất. Các thể của chất
% Don vi : 2.2. Sự chuyển thể của chất (nóng chảy, đông đặc, hoá hơi, ngưng tụ, sôi)
% Muc do : Nhận biết
% YCCD : – Nêu được khái niệm về sự nóng chảy, sự đông đặc, sự bay hơi, sự ngưng tụ và sự sôi.
\begin{ex}%[6K1N2-1]
Quá trình chất chuyển từ thể lỏng sang thể rắn được gọi là
\choice
{sự nóng chảy}
{sự bay hơi}
{\True sự đông đặc}
{sự ngưng tụ}
\loigiai{
Lỏng $\rightarrow$ rắn là sự đông đặc; rắn $\rightarrow$ lỏng là sự nóng chảy; lỏng $\rightarrow$ khí là sự bay hơi; khí $\rightarrow$ lỏng là sự ngưng tụ.
}
\end{ex}

%%%=========EX_11================%%%
% Chu de : Đa dạng chất. Các thể của chất
% Don vi : 2.2. Sự chuyển thể của chất (nóng chảy, đông đặc, hoá hơi, ngưng tụ, sôi)
% Muc do : Nhận biết
% YCCD : – Nêu được nhiệt độ nóng chảy của nước đá và nhiệt độ sôi của nước ở điều kiện thường.
\begin{ex}%[6K1N2-3]
Ở điều kiện thường (áp suất khí quyển chuẩn), nước sôi ở nhiệt độ nào?
\choice
{0 $^\circ$C}
{37 $^\circ$C}
{\True 100 $^\circ$C}
{1000 $^\circ$C}
\loigiai{
Ở áp suất khí quyển chuẩn, nước đá nóng chảy ở 0 $^\circ$C và nước sôi ở 100 $^\circ$C.
}
\end{ex}

%%%=========EX_12================%%%
% Chu de : Đa dạng chất. Các thể của chất
% Don vi : 2.2. Sự chuyển thể của chất (nóng chảy, đông đặc, hoá hơi, ngưng tụ, sôi)
% Muc do : Thông hiểu
% YCCD : – Xác định được quá trình chuyển thể xảy ra trong một hiện tượng cụ thể; tính toán đơn giản với các mốc nhiệt độ chuyển thể.
\begin{ex}%[6K1H2-2]
Sáng sớm, trên lá cây thường có những giọt sương. Hiện tượng này liên quan trực tiếp đến quá trình nào?
\choice
{Sự nóng chảy}
{\True Sự ngưng tụ}
{Sự đông đặc}
{Sự sôi}
\loigiai{
Ban đêm nhiệt độ hạ thấp, hơi nước trong không khí gặp lạnh chuyển thành nước lỏng đọng trên lá -- đó là sự ngưng tụ.
}
\end{ex}
\Closesolutionfile{ans}
\Closesolutionfile{ansex}
%\bangdapan{Ans-KTDK_KHTN6_MoDau_DaDangChat_MADE101}

%%%==============Phần trắc nghiệm đúng sai==============%%%
\subsection{Trắc nghiệm đúng sai}\textit{\large Thí sinh trả lời từ câu 1 đến câu 4. Trong mỗi ý a), b), c), d) ở mỗi câu thí sinh chọn đúng hoặc sai}
\Opensolutionfile{ansex}[Ans/LGTF-KTDK_KHTN6_MoDau_DaDangChat_MADE101]
\Opensolutionfile{ansbook}[Ansbook/AnsTF-KTDK_KHTN6_MoDau_DaDangChat_MADE101]
\Opensolutionfile{ans}[Ans/Tempt-KTDK_KHTN6_MoDau_DaDangChat_MADE101]
\setcounter{ex}{0}
%%%=========TF_1================%%%
% Chu de : Mở đầu về Khoa học tự nhiên
% Don vi : 1.1. Giới thiệu về Khoa học tự nhiên. Các lĩnh vực chủ yếu của KHTN. Vật sống – vật không sống
% Muc do : Nhận biết / Thông hiểu
\begin{ex}%[6KCH1-2]
Xét các nhận định về lĩnh vực của khoa học tự nhiên và về vật sống:
\choiceTF
{\True Vật lí, Hoá học, Sinh học, Thiên văn học và Khoa học Trái Đất là các lĩnh vực chủ yếu của khoa học tự nhiên}
{Nghiên cứu phong tục, tập quán của một dân tộc thuộc lĩnh vực khoa học tự nhiên}
{\True Con gà và cây lúa là vật sống vì chúng có trao đổi chất, lớn lên, sinh sản và cảm ứng}
{Ngọn nến đang cháy được coi là vật sống vì nó toả nhiệt và phát sáng}
\loigiai{
\begin{itemchoice}[T1,F2,T3,F4]
\itemch Đây là 5 lĩnh vực chủ yếu của khoa học tự nhiên.
\itemch Đó là đối tượng của khoa học xã hội và nhân văn.
\itemch Đó là 4 đặc điểm đặc trưng dùng để nhận biết vật sống.
\itemch Toả nhiệt, phát sáng không phải là dấu hiệu của sự sống; ngọn nến không trao đổi chất, không lớn lên và không sinh sản.
\end{itemchoice}
}
\end{ex}

%%%=========TF_2================%%%
% Chu de : Mở đầu về Khoa học tự nhiên
% Don vi : 1.2. An toàn trong phòng thực hành. Sử dụng dụng cụ đo, kính lúp và kính hiển vi quang học
% Muc do : Nhận biết / Thông hiểu
\begin{ex}%[6KCH2-1]
Khi học trong phòng thực hành môn Khoa học tự nhiên, học sinh phải tuân thủ nội quy an toàn và sử dụng đúng dụng cụ đo. Xét các nhận định sau:
\choiceTF
{\True Phải đọc kĩ nhãn và hướng dẫn sử dụng trước khi dùng bất kì hoá chất nào}
{Có thể nếm hoặc ngửi trực tiếp hoá chất để nhận biết chúng}
{\True Khi đun nóng ống nghiệm, phải hướng miệng ống nghiệm về phía không có người}
{Một bình chia độ có GHĐ 100 mL và ĐCNN 2 mL có thể đo chính xác thể tích 55 mL nước}
\loigiai{
\begin{itemchoice}[T1,F2,T3,F4]
\itemch Nhãn hoá chất cho biết tên, nồng độ và các cảnh báo nguy hiểm, bắt buộc phải đọc trước khi sử dụng.
\itemch Tuyệt đối không nếm, không ngửi trực tiếp hoá chất vì có thể gây ngộ độc hoặc bỏng đường hô hấp.
\itemch Chất lỏng sôi có thể bắn ra khỏi miệng ống nghiệm gây bỏng cho người đối diện.
\itemch Với ĐCNN 2 mL, các vạch chia là 0; 2; 4; \dots mL nên chỉ đọc được các giá trị chẵn; 55 mL không trùng vạch chia nào.
\end{itemchoice}
}
\end{ex}

%%%=========TF_3================%%%
% Chu de : Đa dạng chất. Các thể của chất
% Don vi : 2.1. Sự đa dạng của chất. Ba thể cơ bản của chất và đặc điểm
% Muc do : Nhận biết / Thông hiểu / Vận dụng
\begin{ex}%[6K1V1-2]
Cho ba chất ở điều kiện thường: sắt, nước và khí oxygen. Xét các nhận định sau:
\choiceTF
{\True Sắt ở thể rắn nên có hình dạng và thể tích xác định}
{Chất tạo nên vật thể chỉ có ở các vật thể nhân tạo, không có ở vật thể tự nhiên}
{\True Khí oxygen chiếm toàn bộ thể tích bình chứa và dễ bị nén hơn nước}
{\True Rót toàn bộ 200 mL nước từ cốc sang một bình tam giác thì thể tích nước đo được vẫn là 200 mL}
\loigiai{
\begin{itemchoice}[T1,F2,T3,T4]
\itemch Chất rắn có hình dạng và thể tích riêng xác định, rất khó bị nén.
\itemch Mọi vật thể, dù tự nhiên hay nhân tạo, đều được tạo nên từ chất.
\itemch Chất khí không có hình dạng và thể tích xác định, các hạt ở xa nhau nên rất dễ bị nén; chất lỏng khó bị nén hơn nhiều.
\itemch Nước là chất lỏng: hình dạng đổi theo vật chứa nhưng thể tích được bảo toàn.
\end{itemchoice}
}
\end{ex}

%%%=========TF_4================%%%
% Chu de : Đa dạng chất. Các thể của chất
% Don vi : 2.2. Sự chuyển thể của chất (nóng chảy, đông đặc, hoá hơi, ngưng tụ, sôi)
% Muc do : Nhận biết / Vận dụng
\begin{ex}%[6K1V2-5]
Xét các nhận định về sự chuyển thể của chất:
\choiceTF
{\True Sự nóng chảy là quá trình chất chuyển từ thể rắn sang thể lỏng}
{Sự bay hơi chỉ xảy ra khi chất lỏng đạt tới nhiệt độ sôi}
{\True Sự ngưng tụ là quá trình chất chuyển từ thể khí sang thể lỏng}
{Cho nước đá vào cốc nước ngọt, một lúc sau thấy nước bám ngoài thành cốc; đó là do nước trong cốc thấm qua thành cốc ra ngoài}
\loigiai{
\begin{itemchoice}[T1,F2,T3,F4]
\itemch Đây chính là định nghĩa của sự nóng chảy.
\itemch Sự bay hơi xảy ra ở mọi nhiệt độ và chỉ trên bề mặt chất lỏng; sự sôi mới xảy ra ở nhiệt độ xác định và cả trong lòng chất lỏng.
\itemch Ngưng tụ là quá trình ngược với bay hơi.
\itemch Thành cốc thuỷ tinh không cho nước thấm qua. Hơi nước trong không khí gặp thành cốc lạnh đã ngưng tụ thành giọt nước bám bên ngoài.
\end{itemchoice}
}
\end{ex}
\Closesolutionfile{ans}
\Closesolutionfile{ansbook}
\Closesolutionfile{ansex}
%\bangdapanTF{AnsTF-KTDK_KHTN6_MoDau_DaDangChat_MADE101}

%%==============Phần bài tập trả lời ngắn==============%%%
\subsection{Bài tập trả lời ngắn}\textit{\large Thí sinh trả lời từ câu 1 đến câu 2}
\Opensolutionfile{ansex}[Ans/LGSA-KTDK_KHTN6_MoDau_DaDangChat_MADE101]
\Opensolutionfile{ansexh}[Ans/AnsSA-KTDK_KHTN6_MoDau_DaDangChat_MADE101]
\setcounter{ex}{0}
%%%=========SA_1================%%%
% Chu de : Mở đầu về Khoa học tự nhiên
% Don vi : 1.2. An toàn trong phòng thực hành. Sử dụng dụng cụ đo, kính lúp và kính hiển vi quang học
% Muc do : Nhận biết
% YCCD : – Trình bày được công dụng và cách sử dụng một số dụng cụ đo thông dụng (thước, bình chia độ, cân, nhiệt kế), kính lúp và kính hiển vi quang học.
\begin{ex}%[6KCN2-4]
Một kính hiển vi quang học có vật kính ghi 40x và thị kính ghi 10x. Khi dùng cặp vật kính -- thị kính này, ảnh của mẫu vật được phóng đại bao nhiêu lần?
\shortans{$400$}
\loigiai{
Độ phóng đại của kính hiển vi = số bội giác vật kính $\times$ số bội giác thị kính = 40 $\times$ 10 = 400 (lần).
}
\end{ex}

%%%=========SA_2================%%%
% Chu de : Đa dạng chất. Các thể của chất
% Don vi : 2.2. Sự chuyển thể của chất (nóng chảy, đông đặc, hoá hơi, ngưng tụ, sôi)
% Muc do : Thông hiểu
% YCCD : – Xác định được quá trình chuyển thể xảy ra trong một hiện tượng cụ thể; tính toán đơn giản với các mốc nhiệt độ chuyển thể.
\begin{ex}%[6K1H2-3]
Ở điều kiện thường, nước đá nóng chảy ở 0 $^\circ$C và nước sôi ở 100 $^\circ$C. Khoảng nhiệt độ mà nước tồn tại ở thể lỏng rộng bao nhiêu độ C?
\shortans{$100$}
\loigiai{
Nước ở thể lỏng trong khoảng từ 0 $^\circ$C đến 100 $^\circ$C, vậy khoảng nhiệt độ rộng 100 $-$ 0 = 100 ($^\circ$C).
}
\end{ex}
\Closesolutionfile{ansexh}
\Closesolutionfile{ansex}
%\bangdapanSA{AnsSA-KTDK_KHTN6_MoDau_DaDangChat_MADE101}

%%%==============Phần bài tập tự luận==============%%%
\subsection{Bài tập tự luận}\textit{\large Thí sinh trả lời từ bài 1 đến bài 3}
\Opensolutionfile{ansbth}[Ans/LGBT-KTDK_KHTN6_MoDau_DaDangChat_MADE101]
\Opensolutionfile{ansbt}[Ans/AnsBT-KTDK_KHTN6_MoDau_DaDangChat_MADE101]
%%%=========BT_1================%%%
% Chu de : Mở đầu về Khoa học tự nhiên
% Don vi : 1.2. An toàn trong phòng thực hành. Sử dụng dụng cụ đo, kính lúp và kính hiển vi quang học
% Muc do : Thông hiểu
\begin{bt}%[6KCH2-1]
Trong phòng thực hành môn Khoa học tự nhiên có hai quy định sau: (1) Không ăn, uống trong phòng thực hành; (2) Khi đun nóng ống nghiệm phải hướng miệng ống nghiệm về phía không có người.\\
a) Giải thích cơ sở khoa học của quy định (1).\\
b) Giải thích cơ sở khoa học của quy định (2).
\loigiai{
a) Bàn ghế, dụng cụ trong phòng thực hành có thể dính hoá chất độc; thức ăn, đồ uống dễ bị nhiễm hoá chất và đi vào cơ thể qua đường tiêu hoá gây ngộ độc.\\
b) Khi đun, chất lỏng sôi có thể trào hoặc bắn mạnh ra khỏi miệng ống nghiệm; nếu miệng ống hướng về phía người sẽ gây bỏng cho bản thân hoặc bạn bên cạnh.
}
\end{bt}

%%%=========BT_2================%%%
% Chu de : Đa dạng chất. Các thể của chất
% Don vi : 2.1. Sự đa dạng của chất. Ba thể cơ bản của chất và đặc điểm
% Muc do : Vận dụng
\begin{bt}%[6K1V1-3]
Bạn An rót đầy 250 mL nước vào một chiếc cốc hình trụ, sau đó đổ toàn bộ lượng nước này sang một bình tam giác.\\
a) Hình dạng và thể tích của lượng nước thay đổi như thế nào? Giải thích dựa trên đặc điểm của chất ở thể lỏng.\\
b) Nếu thay lượng nước trên bằng một viên nước đá có thể tích 250 cm$^3$ thì hình dạng của nó khi đặt vào bình tam giác có thay đổi không? Vì sao?
\loigiai{
a) Hình dạng thay đổi (từ hình trụ sang hình dạng của bình tam giác) nhưng thể tích không đổi, vẫn là 250 mL. Vì chất ở thể lỏng có thể tích xác định, còn hình dạng thì phụ thuộc vật chứa.\\
b) Không thay đổi. Viên nước đá là chất ở thể rắn, có cả hình dạng và thể tích xác định, không phụ thuộc vào vật chứa.
}
\end{bt}

%%%=========BT_3================%%%
% Chu de : Đa dạng chất. Các thể của chất
% Don vi : 2.2. Sự chuyển thể của chất (nóng chảy, đông đặc, hoá hơi, ngưng tụ, sôi)
% Muc do : Liên hệ thực tế
\begin{bt}%[6K1T2-5]
Vào những ngày trời nồm ở miền Bắc nước ta, sàn nhà lát gạch men và tường nhà thường bị \rq\rq đổ mồ hôi\rq\rq , ướt và trơn trượt rất nguy hiểm.\\
a) Hiện tượng trên liên quan đến quá trình chuyển thể nào của nước? Giải thích vì sao hiện tượng xảy ra mạnh trên nền gạch men và tường nhà.\\
b) Đề xuất 2 biện pháp mà gia đình em có thể thực hiện để hạn chế hiện tượng này và nêu cơ sở khoa học của mỗi biện pháp.
\loigiai{
a) Đó là sự ngưng tụ của hơi nước. Ngày trời nồm, không khí có độ ẩm rất cao (nhiều hơi nước) trong khi nền gạch men và tường nhà còn lạnh do giữ nhiệt của đợt lạnh trước. Hơi nước trong không khí tiếp xúc với bề mặt lạnh, bị làm lạnh xuống dưới điểm sương nên ngưng tụ thành các giọt nước bám trên bề mặt. Gạch men và tường nhẵn, dẫn nhiệt tốt và ấm lên chậm nên hiện tượng xảy ra rõ nhất.\\
b) Nêu đúng mỗi biện pháp kèm cơ sở khoa học được 0,25 điểm, ví dụ: (1) Đóng kín cửa sổ, cửa ra vào để hạn chế không khí ẩm bên ngoài tràn vào nhà, làm giảm lượng hơi nước có thể ngưng tụ. (2) Dùng máy hút ẩm hoặc bật điều hoà ở chế độ khô để giảm độ ẩm không khí trong phòng. (3) Lau sàn bằng khăn khô, giẻ khô thay vì nước, đồng thời rải thảm/giấy báo để chống trơn trượt. (4) Bật quạt, đèn sưởi làm ấm bề mặt sàn -- tường để nhiệt độ bề mặt cao hơn điểm sương thì hơi nước không ngưng tụ.
}
\end{bt}
\Closesolutionfile{ansbt}
\Closesolutionfile{ansbth}

\begin{center}
 \rule[4pt]{2cm}{1pt}\,\large\bfseries Hết\,\rule[4pt]{2cm}{1pt}
\end{center}
\label{x}
\end{document}
